\documentclass[11pt]{article}
\PassOptionsToPackage{obeyspaces,spaces}{url}  % keep spaces in \code and \path

% ============================================================
% COMPILATION
% ============================================================
%
% Required compiler: XeLaTeX
%
% Recommended build sequence:
% XeLaTeX -> Biber -> XeLaTeX -> XeLaTeX
%

% ============================================================
% DOCUMENT EDITION
% ============================================================

\newif\iftutornotes

\tutornotestrue       % TA / Tutor edition
% \tutornotesfalse    % Student edition

% ============================================================
% PACKAGES
% ============================================================

\usepackage[a4paper,margin=1in]{geometry}

% XeLaTeX font system
\usepackage{fontspec}
\usepackage{unicode-math}

\setmainfont{TeX Gyre Termes}
\setsansfont{TeX Gyre Heros}
\setmonofont{Latin Modern Mono}
\setmathfont{TeX Gyre Termes Math}

\usepackage{microtype}
\usepackage{enumitem}
\usepackage{booktabs}
\usepackage{array}
\usepackage{longtable}
\usepackage{xcolor}
\usepackage{hyperref}
\usepackage{listings}
\usepackage{tcolorbox}
\usepackage{titlesec}
\usepackage{fancyhdr}
\usepackage[backend=biber,style=ieee]{biblatex}
\usepackage{xurl}
\usepackage{indentfirst}
\usepackage{graphicx}
%\usepackage{background}
\usepackage{subcaption}

% ============================================================
% DOCUMENT METADATA
% ============================================================

\newcommand{\CourseTitle}{Physical AI, TAICA}
\newcommand{\InstructorName}{Human-centered Intelligent System Lab}
\newcommand{\InstructorEmail}{National Yang Ming Chiao Tung University}

\newcommand{\AssignmentTitle}{Course Project: Robot Learning from Evaluation Feedback}
\newcommand{\Semester}{Fall 2026}

% ============================================================
% BIBLIOGRAPHY
% ============================================================

\addbibresource{refs.bib}

% ============================================================
% TABLE SETTINGS
% ============================================================

\newcolumntype{L}[1]{>{\raggedright\arraybackslash}p{#1}}

% ============================================================
% HYPERLINK SETTINGS
% ============================================================

\hypersetup{
    colorlinks=true,
    linkcolor=blue!50!black,
    citecolor=blue!50!black,
    urlcolor=blue!50!black
}

% ============================================================
% PAGE STYLE
% ============================================================

\setlength{\headheight}{14pt}

\pagestyle{fancy}
\fancyhf{}

\lhead{\CourseTitle}
\rhead{\AssignmentTitle}
\cfoot{\thepage}

% ============================================================
% LIST SETTINGS
% ============================================================

\setlist[itemize]{
    topsep=3pt,
    itemsep=2pt,
    parsep=0pt,
    leftmargin=1.5em
}

\setlist[enumerate]{
    topsep=3pt,
    itemsep=2pt,
    parsep=0pt,
    leftmargin=1.6em
}

% ============================================================
% COLORS
% ============================================================

\definecolor{codebg}{RGB}{248,248,248}
\definecolor{commentgreen}{RGB}{0,110,60}
\definecolor{keywordblue}{RGB}{20,70,150}

% A single restrained color for all internal teaching notes.
\definecolor{tutorblue}{RGB}{65,88,110}
\definecolor{tutorbg}{RGB}{247,249,251}


% ============================================================
% TA / TUTOR NOTE SYSTEM
% ============================================================

% All material intended only for teaching staff must use this
% environment. Do not use separate teacher/grading/TA colors.
%
% In student mode, the complete environment body disappears.

\NewDocumentEnvironment{tutornote}{+b}
{%
    \iftutornotes
        \begin{tcolorbox}[
            colback=tutorbg,
            colframe=tutorblue,
            boxrule=0.4pt,
            arc=1mm,
            left=7pt,
            right=7pt,
            top=6pt,
            bottom=6pt,
            before skip=8pt,
            after skip=8pt
        ]
        {\small
        \color{tutorblue}
        \textbf{TA/Tutor Note.}\par
        \medskip
        #1}
        \end{tcolorbox}
    \fi
}
{}

% Optional edition indicator.
% It appears only in the teaching edition.

\newcommand{\EditionLabel}{%
    \iftutornotes
        \begin{center}
            {\small\color{tutorblue}
            TA / Tutor Teaching Edition}
        \end{center}
        \vspace{0.5em}
    \fi
}

% ============================================================
% LISTINGS
% ============================================================

\lstdefinelanguage{turtle}{
    alsoletter={:@},
    morekeywords={
        @prefix,
        a,
        owl:,
        rdf:,
        rdfs:,
        xsd:
    },
    sensitive=true,
    morecomment=[l]{\#},
    morestring=[b]"
}

\lstset{
    basicstyle=\ttfamily\small,
    backgroundcolor=\color{codebg},
    breaklines=true,
    frame=single,
    framerule=0.2pt,
    rulecolor=\color{black!20},
    keywordstyle=\color{keywordblue}\bfseries,
    commentstyle=\color{commentgreen},
    stringstyle=\color{black!70},
    columns=fullflexible,
    keepspaces=true,
    showstringspaces=false
}

\lstdefinestyle{codesnippet}{
    numbers=none,
    aboveskip=0.75\baselineskip,
    belowskip=0.85\baselineskip,
    frame=single,
    framerule=0.2pt,
    rulecolor=\color{black!20},
    backgroundcolor=\color{codebg},
    basicstyle=\ttfamily\small,
    breaklines=true,
    columns=fullflexible,
    keepspaces=true,
    showstringspaces=false
}

\DeclareUrlCommand\codeurl{\urlstyle{tt}}
\DeclareRobustCommand{\code}{\codeurl}  % breakable inline code

% ============================================================
% TITLE
% ============================================================

\title{\textbf{\AssignmentTitle}}
\author{
    \InstructorName\\
    \InstructorEmail
}
\date{\Semester}

% ============================================================
% DOCUMENT
% ============================================================

\begin{document}

\maketitle
\EditionLabel


%⠀⠀⠀⠀⠀⠀⠀⠀⢀⢔⠾⢋⠷⢃⠠⠒⠈⠀⢀⣀⢂⢠⣲⢦⡪⠝⠀⢠⠎⠀⠀⠀⠀⠀⣠⣴⢪⡃⠜⠋⠉⠀⠀⠀⠀⠀⠀⠀⠀⠀⠀⠄⡀⠖⠁⠀⢀⠀⠀⡀⠀⠀⠀⠀⢀⣤⠦⡻⠂⠀⠀⠀⢼⡆⠀⠀⠁⡔⡀⡸⠀⢠⠃⠀⢸⣐⣷⣏⠉⠁⠉⢻⡄⠀⠀⡱⠑⢆⣨⠟⠊⠉⠀⠀⠈⠁⠀⠀⠀⠀⠀⠀⠀⠀⠀⠀⠀⠀⠀⠀⠀⠀⠀⠀⠀⠀⠀⠀
%⠀⠀⠀⠀⠀⠀⢀⢔⠕⢁⠔⠁⠐⣁⣤⠴⠚⠉⢀⣠⠖⡫⠃⠁⠀⠀⣰⠃⠀⠀⡠⠀⢠⠞⡵⠃⠁⠀⠀⠀⠀⠀⠀⠀⠀⠀⢀⣀⠀⢀⠔⠈⠀⠀⠀⠔⡰⢃⠔⠁⠀⠀⠀⠼⡫⣠⠎⠀⠀⢀⠤⠀⠀⠀⠀⠀⠀⠰⢁⠃⢀⠇⠀⠀⣼⠋⣟⡆⠇⢀⠀⠸⢎⢵⠀⠱⠱⡈⢧⠄⠀⠀⠀⠀⠀⠀⠀⠀⠀⠀⠀⠀⠀⠀⠀⠀⠀⠀⠀⠀⠀⠀⠀⠀⠀⠀⠀⠀
%⠀⠀⠀⠀⠀⣠⡡⠁⢠⣁⣀⡴⡚⠅⠐⠈⢀⡴⠋⠐⠁⢀⡠⡤⠄⣰⠃⠀⠀⣰⠁⠀⣠⠊⠀⠀⠀⠀⢀⣀⣠⣤⠤⠶⠚⠋⠉⢀⡠⠀⠀⠀⠀⢠⣮⠞⣐⠅⠀⠀⠀⠀⢀⡴⠋⠁⣀⣠⠔⠁⠀⠀⢠⠃⠀⠀⠀⢦⠂⢠⠊⠀⠀⠚⠙⢰⢸⣷⢰⠈⢆⠀⠙⣮⢣⠀⠐⡔⡒⣧⠀⠀⠀⠀⠀⠀⠀⠀⠀⠀⠀⠀⠀⠀⠀⠀⠀⠀⠀⠀⠀⠀⠀⠀⠀⠀⠀⠀
%⠀⠀⠀⠀⢠⠟⣠⠖⠋⢁⠚⠃⣀⠔⠚⡷⠉⣉⡤⡲⢭⠞⢉⠃⣰⠏⠀⠀⡴⣉⡀⡚⠁⠑⠒⠀⡛⠛⢉⢁⠄⠀⣠⠗⠀⢀⡴⠋⠀⢀⠤⢠⢾⠋⢡⠞⠁⠀⡠⠒⠀⢀⣎⣠⢞⢵⠟⠁⠀⢀⠔⠠⠃⠀⡔⡐⠀⡇⠀⠁⠀⠀⠀⠀⡃⠈⠀⣿⠈⡀⡇⠣⡀⠈⢧⠡⡀⠈⢊⢜⣧⡀⠀⠀⠀⠀⠀⠀⠀⣀⡠⢤⣤⡀⠀⠀⠀⠀⠀⠀⠀⠀⠀⠀⠀⠀⠀⠀
%⠀⠀⠀⠀⢸⢶⠃⠀⢀⠃⢠⡞⠁⢀⡼⡷⢋⣥⣮⠴⠁⡠⣵⢻⡟⠀⢀⡼⢋⢊⠌⠀⡠⠊⢀⠊⢀⣀⣆⠃⠀⣰⠃⢀⡴⠋⠀⢀⠔⠕⡡⠞⣠⠝⠁⠀⣠⠊⢀⣤⠖⣡⠞⣕⡡⠁⠀⣠⡞⡡⣶⡵⠀⣸⢣⠇⢸⡟⠀⠀⠀⡀⠀⢠⠁⡆⠀⡷⠀⡇⣏⡄⢻⠄⠀⠱⡷⣄⠀⠡⡹⡇⠀⢀⡀⠄⠒⠈⠁⠀⠀⠀⠀⠉⠓⢤⡀⠀⠀⠀⠀⠀⠀⠀⠀⠀⠀⠀
%⠀⠀⠀⠀⠈⠙⠀⠀⡎⢠⠋⠀⣠⡮⠔⠈⣩⠞⠁⢀⢊⡾⢡⡄⠁⢠⡾⠡⠡⢂⠠⠊⢀⠔⣀⡴⢋⡏⠎⠀⣸⠃⣰⠟⠁⠀⡐⠁⡡⡊⠔⠈⡁⢐⣔⡟⢡⠞⡑⣡⠎⣡⠞⠝⠀⢀⣮⢟⠊⡸⠹⠁⢰⠃⣼⠀⣿⠂⢰⠀⢰⠃⠀⡌⢰⡗⢰⡧⠀⢫⡷⠇⠀⣎⢆⠀⡇⠏⢳⡤⠜⠓⠈⠁⠀⠀⠀⠀⠀⠀⠀⠀⠀⠀⠀⠀⠈⠓⢄⡀⠀⠀⠀⠀⠀⠀⠀⠀
%⠀⠀⠀⠀⠀⠀⠀⠀⠑⣿⣖⡾⠋⠀⡠⠊⠁⢠⣖⣵⡭⡂⠁⠘⠄⢳⠁⡶⠓⣡⣰⣖⡥⠞⠁⣀⢼⢱⠀⣰⢃⡼⠃⠀⢠⡪⣪⠞⠋⡀⢔⣠⠦⠛⠉⣠⡳⢊⡴⢣⠞⢁⠊⠀⣠⡿⠛⢁⠎⢠⡳⢡⠏⢸⠟⢠⢟⠀⢸⠀⢸⠁⠀⢁⡿⠁⢠⠇⠀⠘⡇⠘⠀⡆⣾⠠⠓⠉⠁⠀⠀⠀⠀⠀⠀⠀⠀⠀⠀⠀⠀⠀⠀⠀⠀⠀⠀⠀⠀⠙⠦⡀⠀⠀⠀⠀⠀⠀
%⠀⠀⠀⠀⠀⠀⠀⠀⡀⠝⠛⠀⡠⢊⣠⠾⠗⡾⠁⢳⠀⣽⡄⠀⢘⠾⣊⠴⢋⡵⢫⣷⣃⢀⠔⠁⣿⠆⠀⠀⡞⢁⠀⠴⠛⠘⣀⣔⡬⢖⠋⠁⢀⠔⣶⡟⠡⢈⡕⠛⠠⠂⠀⣰⠋⠀⡰⠁⢀⣧⢡⡎⢀⠟⠀⣸⢹⠀⢸⠀⡸⠀⠀⣾⠁⠀⡌⠀⣀⡸⠙⣄⠶⠋⠀⠀⠀⠀⠀⠀⠀⠀⠀⠀⠀⠀⠀⠀⠀⠀⠀⠀⠀⠀⠀⠀⠀⠀⠀⠀⠀⠉⠦⡀⠀⠀⠀⠀
%⠀⠀⠀⠀⠀⠀⠀⠉⠀⠀⣰⣪⠖⠋⠀⢀⠜⢀⢌⠜⡆⣿⠼⣺⢗⣟⣡⡎⢡⠃⣤⠹⠘⠢⡤⢄⣛⡐⠠⠼⠍⠐⠀⣀⡤⡞⠉⠁⣤⠋⢀⠔⠁⣼⠏⠐⢠⠎⠐⠰⠃⠀⡼⠁⠀⡜⠀⡰⣻⢃⡞⠀⣸⠃⢀⡟⢰⠀⢸⠀⣷⠀⣸⠁⠀⡘⠀⣼⡿⠁⢠⢏⠀⠀⠀⠀⠀⠀⠀⠀⠀⠀⠀⠀⠀⠀⠀⠀⠀⠀⠀⠀⠀⠀⠀⠀⠀⠀⠀⠀⠀⠀⠀⠙⢦⠀⠀⠀
%⠀⠀⠀⠀⠀⠀⠀⠀⠀⢎⣸⠃⠀⠀⡠⢃⣴⠟⣡⣴⣿⣷⠛⠛⡈⡇⢠⠗⢸⣾⠸⡇⠀⠀⠈⠑⡾⣫⢒⣴⣶⢟⠵⢡⠌⠀⠀⠔⠃⡠⠁⠀⣾⢋⠌⡰⠁⠀⠠⢁⡄⠐⠀⠀⢞⡒⡰⢠⡏⡾⢠⢧⡏⠀⢺⠯⢥⠀⢸⠀⡽⢠⠃⠀⡰⠇⣸⠗⠃⢠⢳⡆⠀⠀⠀⠀⠀⠀⠀⠀⠀⠀⠀⠀⠀⠀⠀⠀⠀⠀⠀⠀⠀⠀⠀⠀⠀⠀⠀⠀⠀⠀⠀⠀⠈⠷⡄⠀
%⠀⠀⠀⠀⠀⠀⠀⠀⠀⠀⣨⣢⢔⣡⠖⢫⠽⠑⡟⠉⢸⢯⢏⠉⣵⡇⣸⠀⢘⢨⠃⡁⠓⠒⢢⢞⢜⣥⣫⣿⡧⠃⠀⡌⠀⠀⠈⠀⠊⠀⠀⣼⠃⠊⡐⠀⠀⢠⠣⡞⢠⡆⠀⡎⠀⠀⠁⣼⡝⢠⠟⡸⠀⠀⢸⢐⣸⠂⢸⠀⡇⢂⠄⡠⣧⣷⠏⠀⢠⡇⡈⡇⠀⠀⠀⠀⠀⠀⠀⠀⠀⣴⡆⠀⢰⡆⠀⢠⣾⣷⠀⢀⣾⠀⠀⣾⠆⠀⠀⠀⠀⠀⠀⠀⠀⠀⢹⡆
%⠀⠀⠀⠀⠀⠀⠀⠀⠀⠚⠿⠚⠋⠁⠌⠁⡠⠊⡐⣡⢟⡌⡈⢒⡏⢰⢸⢰⢸⢺⡄⠀⢀⡴⣷⢿⠏⢈⣿⠟⠷⢆⡤⢀⣀⠀⠀⠀⠀⠀⢀⣏⠌⡔⠀⠀⠀⣆⠾⢁⡞⡇⡜⠀⠀⠀⣀⣯⡴⣥⢷⠓⠒⠋⠉⠡⢸⠀⢸⠀⠇⠎⣠⠱⢸⠇⠀⢠⠃⢠⢁⡇⠀⠀⠀⠀⠀⠀⠀⠀⢰⣿⣷⢀⡿⢁⣴⠟⣹⣿⠀⣾⣷⣶⣾⡟⠀⠀⠀⠀⠀⠀⠀⠀⠀⠀⠸⡇
%⠀⠀⠀⠀⠀⠀⠀⠀⠀⠀⠀⠀⠀⣠⠔⠋⣀⣪⠞⡑⢁⢧⠙⠀⢇⢸⣿⢺⡞⡚⡯⢴⡙⠉⢀⡐⠂⡘⡞⠀⠅⠠⠉⠁⠚⠣⠝⣔⠶⣀⠀⠁⡰⠀⠀⠀⠀⡏⣄⡜⠀⡷⢓⣢⠿⠍⠛⠋⠠⡡⠂⠀⠀⠀⢀⠀⢀⡇⢨⠀⣶⡜⢸⠀⡟⠀⠀⢆⣠⠃⡈⡇⠀⠀⠀⠀⠀⠀⠀⢀⣿⠉⣿⣿⢃⣾⠿⠟⢻⣟⣸⡟⠁⢠⡿⠁⢀⣴⠆⠀⠀⠀⠀⠀⠀⠀⢠⡇
%⠀⠀⠀⠀⠀⠀⠀⠀⠀⠀⣀⠴⠞⠧⠜⠋⢁⠂⠌⡰⢃⣾⣦⠀⢸⡀⣿⣄⢁⠇⠀⠀⢯⣝⣖⣿⣯⣿⣒⡭⠥⣐⡒⠤⢀⠀⠀⠈⠉⡛⢆⢠⠃⠀⠀⠀⠀⠑⠋⢀⣀⠻⠉⠁⠀⠀⢀⡠⠔⣒⣀⣭⣝⣛⣫⣿⣿⣧⠘⠀⣳⠀⣼⢠⠃⠀⠀⠸⣹⢀⠃⡇⠀⠀⠀⠀⠀⠀⠀⠘⠃⠀⠘⠃⠘⠃⠀⠀⠙⠋⠛⠀⠀⠙⠃⢠⠞⠁⠀⠀⠀⠀⠀⠀⠀⠀⢸⠃
%⠀⠀⠀⠀⠀⠀⠀⠀⠀⠈⠁⠀⠀⠀⠀⢀⠂⣼⡞⠁⣾⡡⠘⡆⢸⡇⢸⡘⡈⢶⠀⠀⠀⢼⣣⠀⠀⢸⡄⠈⢽⡇⠛⠷⣦⣽⠀⠀⠀⠘⡟⡇⠀⠀⠀⠀⠀⠀⢠⣾⠆⠀⠀⠀⠀⢀⣯⡴⣾⠛⠩⠧⠄⣸⠃⠀⠀⣼⠆⢰⣇⡼⢰⣿⠄⠀⠀⠠⢿⠸⢰⡇⠀⠀⠀⠀⠀⠀⠀⠀⠀⠀⠀⢀⣀⣀⣤⡀⣤⡄⢀⣠⣤⣤⣤⣀⠀⠀⠀⠀⠀⠀⠀⠀⠀⠀⢸⠀
%⠀⠀⠀⠀⠀⠀⠀⠀⠀⠀⠀⠀⠀⠀⠀⣨⡴⠋⠀⢰⣿⠁⢀⣿⣸⠁⢶⠧⠀⠘⢇⠀⠀⠘⢞⡄⠀⠀⠙⠧⣀⡂⠤⠂⠈⢿⣇⠀⠀⠀⠼⢿⠀⠀⠀⠀⠀⠀⢠⠘⠆⠀⠀⠀⠀⣼⡏⠀⠈⠢⢄⡤⠴⠋⠀⠀⡰⠸⡄⢳⣷⢇⡂⣯⠀⠀⢠⠃⠎⢀⠟⡇⠀⠀⠀⠀⠀⠀⠀⠀⠀⠀⠀⠙⢻⡟⠉⣿⠋⠀⣿⠟⠉⠉⣽⡟⠀⠀⠀⠀⠀⠀⠀⠀⠀⠀⢸⠀
%⠀⠀⠀⠀⠀⠀⠀⠀⠀⠀⠀⠀⠀⠀⠀⠉⠀⠀⠀⣿⡏⠀⣼⢣⣿⡄⢸⡇⠀⠀⠈⠂⠀⠀⠀⠜⢄⡀⠀⠀⠀⠀⠀⠀⢀⡟⠯⡶⡶⢀⡶⠋⠀⠀⠀⠀⠀⠀⠀⠁⠀⢀⣄⡀⢰⡟⠳⡀⠀⠀⠀⠀⠀⠀⢀⠔⠣⠀⡌⡟⡆⢸⠇⠈⠁⡤⢸⡀⢠⠎⠀⡇⠀⠀⠀⠀⠀⠀⠀⠀⠀⠀⠀⢠⡿⠁⠀⠀⠀⢰⡟⣀⣤⡾⠋⠀⠀⠀⠀⠀⠀⠀⠀⠀⠀⠀⢸⠀
%⠀⠀⠀⠀⠀⠀⠀⠀⠀⠀⠀⠀⠀⠀⠀⠀⠀⠀⠀⡟⢣⠜⠁⡾⠼⡆⠀⣷⠀⠀⠀⠀⠀⠀⠀⠀⠀⠀⠀⠀⠀⠉⠉⠉⠁⠀⠚⠋⠀⠌⠀⠀⠀⠀⠀⠀⠀⠀⠀⠀⠀⠀⠓⣞⣤⡀⠁⠉⠉⠉⠉⠀⠀⠀⠀⠀⠀⠀⣷⢸⡁⠈⠀⠀⢀⡇⠀⢡⢿⠁⠀⡇⠀⠀⠀⠀⠀⠀⠀⠀⠀⠠⣶⡾⠿⠟⠂⠀⠀⠿⠿⠛⠉⠀⠀⠀⠀⠀⠀⠀⠀⠀⠀⠀⠀⠀⢸⠀
%⠀⠀⠀⠀⠀⠀⠀⠀⠀⠀⠀⠀⠀⠀⠀⠀⠀⠀⠄⠊⠁⠀⡸⡷⢣⠺⣆⢸⡇⠀⠀⠀⠀⠀⠀⠀⠀⠀⠀⠀⠀⠀⠀⠀⠀⠀⠀⠀⠁⠀⠀⠀⠀⠀⠀⠀⠀⠀⠀⠀⠀⠀⠀⠀⠀⠀⠀⠀⠀⠀⠀⠀⠀⠀⠀⠀⠀⡰⠇⠘⡄⠀⠀⠀⢸⠀⢠⢿⢹⠀⠀⡇⠀⠀⠀⠀⠀⠀⠀⠀⠀⠀⠀⠀⠀⠀⠀⠀⡀⠀⠀⠀⠀⠀⠀⢀⠀⠀⠀⠀⠀⠀⠀⠀⠀⠀⢸⠀
%⠀⠀⠀⠀⠀⠀⠀⠀⠀⠀⠀⠀⠀⠀⠀⠀⠀⠀⠀⠀⠀⣠⡵⠓⠁⠀⠈⠛⣷⠀⠀⠀⠀⠀⠀⠀⠀⠀⠀⠀⠀⠀⠀⠀⠀⠀⠀⠀⠀⠀⠀⠀⠀⠀⠀⠀⠀⠀⠀⠀⠀⠀⠀⠀⠀⠀⠀⠀⠀⠀⠀⠀⠀⠀⠀⠀⠐⠁⠀⠀⠁⠀⠀⢀⣇⣠⣟⡟⡜⡀⠀⡇⠀⠀⠀⠀⠀⠀⠀⠀⠀⣾⡆⢠⣷⡆⠀⣼⠇⣰⡇⢰⣷⠀⢀⣾⠃⠀⠀⠀⠀⠀⠀⠀⠀⠀⢸⠀
%⠀⠀⠀⠀⠀⠀⠀⠀⠀⠀⠀⠀⠀⠀⠀⠀⠀⠀⠀⠀⠀⠁⠀⠀⠀⠀⠀⠀⢹⡆⠀⠀⠀⠀⠀⠀⠀⠀⠀⠀⠀⠀⠀⠀⠀⠀⠀⠀⠀⠀⠀⠀⠀⠀⠀⠀⠀⠀⠀⠀⠀⠀⠀⠀⠀⠀⠀⠀⠀⠀⠀⠀⠀⠀⠀⠀⠀⠀⠀⠀⠀⠀⠀⣼⠋⣿⢰⠇⢸⣱⡀⢤⠀⠀⠀⠀⠀⠀⠀⠀⢸⣿⢠⡾⣿⣇⣼⠏⢠⡿⢀⣿⢿⣧⣼⠇⠀⠀⠀⠀⠀⠀⠀⠀⠀⠀⢸⠂
%⠀⠀⠀⠀⠀⠀⠀⠀⠀⠀⠀⠀⠀⠀⠀⠀⠀⠀⠀⠀⠀⠀⠀⠀⠀⠀⠀⠀⠘⣧⠀⠀⠀⠀⠀⠀⠀⠀⠀⠀⠀⠀⠀⠀⠀⠀⠀⠀⠀⠀⠀⠀⠀⠀⠀⠀⠀⠀⠀⠀⠀⠀⠀⠀⠀⠀⠀⠀⠀⠀⠀⠀⠀⠀⠀⠀⠀⠀⠀⠀⠀⠀⣼⠇⠀⣿⣿⠀⠀⡏⠁⢸⠀⠀⠀⠀⠀⠀⠀⠀⢸⣿⡟⠁⢿⣿⠋⢀⣾⠁⣾⠃⠈⣿⡟⠀⣤⡦⠀⠀⠀⠀⠀⠀⠀⠀⢸⠀
%⠀⠀⠀⠀⠀⠀⠀⠀⠀⠀⠀⢀⡀⠀⠀⠀⠀⠀⠀⠀⠀⠀⠀⠀⠀⠀⠀⠀⠀⢹⡄⠀⠀⠀⠀⠀⠀⠀⠀⠀⠀⠀⠀⠀⠀⠀⠀⠀⠀⠀⠀⠀⠀⠀⠀⠀⠀⠀⠀⠀⠀⠀⠀⠀⠀⠀⠀⠀⠀⠀⠀⠀⠀⠀⠀⠀⠀⠀⠀⠀⠀⢰⡟⠀⢰⣟⠇⠀⠈⠀⠀⡼⠗⠀⠀⠀⠀⠀⠀⠀⠈⠉⠀⠀⠈⠁⠀⠈⠃⠀⠉⠀⠀⠈⠁⠀⠉⠀⠀⠀⠀⠀⠀⠀⠀⢀⠞⠀
%⠀⠀⠀⠀⠀⠀⠀⠀⠀⠀⠀⠀⠀⠀⠀⠀⠀⠀⠀⠀⠀⠀⠀⠀⠀⠀⠀⠀⠀⠈⣇⠀⠀⠀⠀⠀⠀⠀⠀⠀⠀⠀⠀⠀⠀⠀⠀⠀⠀⠀⠀⠀⠀⠀⠀⠀⠀⠀⠀⠀⠀⠀⠀⠀⠀⠀⠀⠀⠀⠀⠀⠀⠀⠀⠀⠀⠀⠀⠀⠀⢀⣿⠁⠀⡾⠊⠀⠀⠀⠀⠀⠹⣄⠀⠀⠀⠀⠀⠀⠀⠀⠀⠀⠀⠀⠀⠀⠀⠀⠀⠀⠀⠀⠀⠀⠀⠀⠀⠀⠀⠀⠀⠀⠀⢀⠎⠀⠀
%⠀⠀⠀⠀⠀⠀⠀⠀⠀⠀⠀⠀⠀⠀⠀⠀⠀⠀⠀⠀⠀⠀⠀⠀⠀⠀⠀⠀⠀⠀⢻⡀⠀⠀⠀⠀⠀⠀⠀⠀⠀⠀⠀⠀⠀⠀⠀⠀⠀⠀⠀⣀⠀⠀⠀⠀⠀⠀⠀⠀⠀⠀⠀⠀⠀⠀⠀⠀⠀⠀⠀⠀⠀⠀⠀⠀⠀⠀⠀⠀⣼⠇⠀⠀⠀⠀⠀⠀⠀⠀⠀⠀⠈⠓⢄⠀⠀⠀⠀⠀⠀⠀⠀⠀⠀⠀⠀⠀⠀⠀⠀⠀⠀⠀⠀⠀⠀⠀⠀⠀⠀⠀⠀⣠⠋⠀⠀⠀
%⠀⠀⠀⠀⠀⠀⠀⠀⠀⠀⠀⠀⠀⠀⠀⠀⠀⠀⠀⠀⠀⠀⠀⠀⠀⠀⠀⠀⠀⠀⠈⣷⡄⠀⠀⠀⠀⠀⠀⠀⠀⠀⠀⠀⠀⠀⠀⠀⠀⠀⠀⠀⠋⢦⡀⠀⢀⠄⠀⠀⠀⠀⠀⠀⠀⠀⠀⠀⠀⠀⠀⠀⠀⠀⠀⠀⠀⠀⠀⣴⡟⠀⠀⠀⠀⠀⠀⠀⠀⠀⠀⠀⠀⠀⠀⠈⠂⢄⡀⠀⠀⠀⠀⠀⠀⠀⠀⠀⠀⠀⠀⠀⠀⠀⠀⠀⠀⠀⠀⠀⠀⠀⣰⠃⠀⠀⠀⠀
%⠀⠀⠀⠀⠀⠀⠀⠀⠀⠀⠀⠀⠀⠀⠀⠀⠀⠀⠀⠀⠀⠀⠀⠀⠀⠀⠀⠀⠀⠀⠀⢻⡽⢆⠀⠀⠀⠀⠀⠀⠀⠀⠀⠀⠀⠀⠀⠀⠀⠀⠀⠀⠀⠈⠁⠀⠈⠀⠀⠀⠀⠀⠀⠀⠀⠀⠀⠀⠀⠀⠀⠀⠀⠀⠀⠀⠀⣠⠞⣿⠇⠀⠀⠀⠀⠀⠀⠀⠀⠀⠀⠀⠀⠀⠀⠀⠀⠀⠈⠐⠢⢄⣀⠀⠀⠀⠀⠀⠀⠀⠀⠀⠀⠀⠀⠀⠀⢀⣀⣠⣴⡾⠁⠀⠀⠀⠀⠀
%⠀⠀⠀⠀⠀⠀⠀⠀⠀⠀⠀⠀⠀⠀⠀⠀⠀⠀⠀⠀⠀⠀⠀⠀⠀⠀⠀⠀⠀⠀⠀⢸⣿⡎⢳⡀⠀⠀⠀⠀⠀⠀⠀⠀⠀⠀⠀⠀⠀⠀⠀⠀⠀⠀⠀⠀⠀⠀⠀⠀⠀⠀⠀⠀⠀⠀⠀⠀⠀⠀⠀⠀⠀⠀⠀⣠⠞⠁⣼⣿⠀⠀⠀⠀⠀⠀⠀⠀⠀⠀⠀⠀⠀⠀⢀⡀⠤⠀⠀⠀⠀⠀⠈⠉⠓⢾⠗⣀⠤⠤⠤⠄⠀⠀⠀⠀⠈⠁⠀⠀⠀⠀⠀⠀⠀⠀⠀⠀
%⠀⠀⠀⠀⠀⠀⠀⠀⠀⠀⠀⠀⠀⠀⠀⠀⠀⠀⠀⠀⠀⠀⠀⠀⠀⠀⠀⠀⠀⠀⠀⢸⣯⣳⠀⠹⣆⠀⠀⠀⠀⠀⠰⠦⠄⣀⣀⡀⠀⠀⠀⠀⠀⢀⡀⠀⠀⠀⠀⠀⠀⠀⠀⠀⣀⡀⠠⠴⠛⠁⠀⠀⠀⢀⡜⠁⠀⡰⣿⡇⠀⠀⠀⠀⠀⠀⠀⠀⠀⠀⠀⠀⠀⠀⠁⠀⠀⠀⠀⠀⠀⠀⠀⠀⢠⢋⠖⠁⠀⠀⠀⠀⠀⠀⠀⠀⠀⠀⠀⠀⠀⠀⠀⠀⠀⠀⠀⠀
%⠀⠀⠀⠀⠀⠀⠀⠀⠀⠀⠀⠀⠀⠀⠀⠀⠀⠀⠀⠀⠀⠀⠀⠀⠀⠀⠀⠀⠀⠀⠀⢸⣿⢌⡆⠀⠈⢷⡀⠀⠀⠀⠀⠀⠀⠀⠀⠨⡙⢿⣿⣛⠛⠛⠛⠛⠛⠛⠻⣛⣩⢿⡋⠉⠀⠀⠀⠀⠀⠀⠀⢀⡴⠋⠀⢀⣴⠿⣿⠅⠀⠀⠀⠀⠀⠀⠀⠀⠀⠀⠀⠀⠀⠀⠀⠀⠀⠀⠀⠀⠀⠀⠀⢠⠗⠁⠀⠀⠀⠀⠀⠀⠀⠀⠀⠀⠀⠀⠀⠀⠀⠀⠀⠀⠀⠀⠀⠀
%⠀⠀⠀⠀⠀⠀⠀⠀⠀⠀⠀⠀⠀⠀⠀⠀⠀⠀⠀⠀⠀⠀⠀⠀⠀⠀⠀⠀⠀⠀⠀⢸⡏⠳⣵⠀⠀⠀⠹⣦⠀⠀⠀⠀⠀⠀⠀⠀⠀⠑⠦⣤⣤⣭⣭⣭⣭⣭⢩⡭⠖⠃⠀⠀⠀⠀⠀⠀⠀⢀⡴⠋⠀⠀⠀⣲⠯⢸⡿⠁⠀⠀⠀⠀⠀⠀⠀⠀⠀⠀⠀⠀⠀⠀⠀⠀⠀⠀⠀⠀⠀⠀⠠⠋⠀⠀⠀⠀⠀⠀⠀⠀⠀⠀⠀⠀⠀⠀⠀⠀⠀⠀⠀⠀⠀⠀⠀⠀
%⠀⠀⠀⠀⠀⠀⠀⠀⠀⠀⠀⠀⠀⠀⠀⠀⠀⠀⠀⠀⠀⠀⠀⠀⠀⠀⠀⠀⠀⠀⠀⢸⡇⠀⠈⠀⠀⠀⠀⠈⢳⡄⠀⠀⠀⠀⠀⠀⠀⠀⠀⠀⠀⠀⠀⠀⠀⠀⠀⠀⠀⠀⠀⠀⠀⠀⠀⠀⡴⠋⠀⠀⠀⠀⢰⡿⢌⣿⡇⠀⠀⠀⠀⠀⠀⠀⠀⠀⠀⠀⠀⠀⠀⠀⠀⠀⠀⠀⠀⠀⠀⠀⠀⠀⠀⠀⠀⠀⠀⠀⠀⠀⠀⠀⠀⠀⠀⠀⠀⠀⠀⠀⠀⠀⠀⠀⠀⠀
%⠀⠀⠀⠀⠀⠀⠀⠀⠀⠀⠀⠀⠀⠀⠀⠀⠀⠀⠀⠀⠀⠀⠀⠀⠀⠀⠀⠀⠀⠀⠀⢸⡇⠀⠀⠀⠀⠀⠀⠀⠩⡿⣦⡀⠀⠀⠀⠀⠀⠀⠀⠀⠀⠀⠀⠀⠀⠀⠀⠀⠀⠀⠀⠀⠀⠀⣠⠎⠀⠀⠀⠀⠀⠀⢐⡺⢁⣿⠃⠀⠀⠀⠀⠀⠀⠀⠀⠀⠀⠀⠀⠀⠀⠀⠀⠀⠀⠀⠀⠀⠀⠀⠀⠀⠀⠀⠀⠀⠀⠀⠀⠀⠀⠀⠀⠀⠀⠀⠀⠀⠀⠀⠀⠀⠀⠀⠀⠀
%⠀⠀⠀⠀⠀⠀⠀⠀⠀⠀⠀⠀⠀⠀⠀⠀⠀⠀⠀⠀⠀⠀⠀⠀⠀⠀⠀⠀⠀⠀⣀⣼⣧⠀⢈⠢⡀⠀⠀⠀⠀⢘⣝⢿⣆⠀⠀⠀⠀⠀⠀⠀⠀⠀⠀⠀⠀⠀⠀⠀⠀⠀⠀⠀⣠⣾⠏⠀⠀⠀⠀⠀⠀⠒⠁⠀⢸⢹⠀⠀⠀⠀⠀⠀⠀⠀⠀⠀⠀⠀⠀⠀⠀⠀⠀⠀⠀⠀⠀⠀⠀⠀⠀⠀⠀⠀⠀⠀⠀⠀⠀⠀⠀⠀⠀⠀⠀⠀⠀⠀⠀⠀⠀⠀⠀⠀⠀⠀
%⠀⠀⠀⠀⠀⠀⠀⠀⠀⠀⠀⠀⠀⠀⠀⠀⠀⠀⠀⠀⠀⠀⠀⠀⠀⠀⣀⣤⣶⣿⣿⣿⣿⠀⠀⠑⢌⢲⣀⠀⠀⠀⠳⣅⡫⠳⣤⡀⠀⠀⠀⠀⠀⠀⠀⠀⠀⠀⠀⠀⠀⠀⣠⣾⡽⠏⠀⠀⠀⠀⠀⠀⠀⠀⠀⠀⣿⡇⠀⠀⠀⠀⠀⠀⠀⠀⠀⠀⠀⠀⠀⠀⠀⠀⠀⠀⠀⠀⠀⠀⠀⠀⠀⠀⠀⠀⠀⠀⠀⠀⠀⠀⠀⠀⠀⠀⠀⠀⠀⠀⠀⠀⠀⠀⠀⠀⠀⠀
%⠀⠀⠀⠀⠀⠀⠀⠀⠀⠀⠀⠀⠀⠀⠀⠀⠀⠀⠀⠀⠀⠀⢀⣠⣴⣾⣿⣿⣿⣿⣿⣿⣷⠀⠀⠀⢳⣍⠣⡀⠀⠀⠀⢻⡟⠆⠈⠹⠶⣤⣀⣀⣀⣀⡀⣀⢀⣀⣀⣠⣴⠿⣛⢝⡃⠀⠀⠀⠀⠀⠀⠀⠀⠀⠀⠀⣿⣿⣶⣤⡀⠀⠀⠀⠀⠀⠀⠀⠀⠀⠀⠀⠀⠀⠀⠀⠀⠀⠀⠀⠀⠀⠀⠀⠀⠀⠀⠀⠀⠀⠀⠀⠀⠀⠀⠀⠀⠀⠀⠀⠀⠀⠀⠀⠀⠀⠀⠀
%⠀⠀⠀⠀⠀⠀⠀⠀⠀⠀⠀⠀⠀⠀⠀⠀⠀⠀⠀⣀⣴⣾⣿⣿⣿⣿⣿⣿⣿⣿⣿⣿⣿⣆⠀⠀⠈⢯⡫⡳⡀⠀⠀⠀⠻⣯⡀⠀⠀⠀⠀⠈⠉⠉⠙⠙⠋⠛⠙⢫⣍⡳⢎⠋⠀⠀⠀⣠⡴⠖⠀⠀⠀⠀⠀⢀⣿⣿⣿⣿⣿⣷⣤⣀⠀⠀⠀⠀⠀⠀⠀⠀⠀⠀⠀⠀⠀⠀⠀⠀⠀⠀⠀⠀⠀⠀⠀⠀⠀⠀⠀⠀⠀⠀⠀⠀⠀⠀⠀⠀⠀⠀⠀⠀⠀⠀⠀⠀
%⠀⠀⠀⠀⠀⠀⠀⠀⠀⠀⠀⠀⠀⠀⢀⣠⣴⣾⣿⣿⣿⣿⣿⣿⣿⣿⣿⣿⣿⣿⣿⣿⣿⣿⣷⣄⠀⠈⡟⣬⡑⠄⠀⠀⠀⠱⡷⣂⠤⣀⠀⠀⠀⠀⠀⠀⠀⠀⠀⣜⣪⡍⠀⠀⠀⠀⡰⠏⠁⣀⠀⠀⠀⠀⠀⣸⣿⣿⣿⣿⣿⣿⣿⣿⣿⣦⣀⠀⠀⠀⠀⠀⠀⠀⠀⠀⠀⠀⠀⠀⠀⠀⠀⠀⠀⠀⠀⠀⠀⠀⠀⠀⠀⠀⠀⠀⠀⠀⠀⠀⠀⠀⠀⠀⠀⠀⠀⠀
%⠀⠀⠀⠀⠀⠀⠀⠀⠀⠀⠀⣀⣤⣾⣿⣿⣿⣿⣿⣿⣿⣿⣿⣿⣿⣿⣿⣿⣿⣿⣿⣿⣿⣿⣿⣿⣶⣄⠸⡦⡋⡦⠀⠀⠀⠀⠈⢎⠙⠕⣓⠂⠤⢀⢠⠄⠀⠀⠜⠳⠋⠀⠀⠀⢀⡾⢋⡥⠴⣇⠀⠀⠀⠀⠀⣿⣿⣿⣿⣿⣿⣿⣿⣿⣿⣿⣿⣿⣶⣤⣀⠀⠀⠀⠀⠀⠀⠀⠀⠀⠀⠀⠀⠀⠀⠀⠀⠀⠀⠀⠀⠀⠀⠀⠀⠀⠀⠀⠀⠀⠀⠀⠀⠀⠀⠀⠀⠀
%⠀⠀⠀⠀⠀⢀⣀⣤⣴⣶⣿⣿⣿⣿⣿⣿⣿⣿⣿⣿⣿⣿⣿⣿⣿⣿⣿⣿⣿⣿⣿⣿⣿⣿⣿⣿⣿⣿⣷⣌⡢⡹⡄⠀⠀⠀⠀⠘⡖⠔⣒⡠⠄⡭⠃⠀⠀⠐⡺⠁⠀⠀⢀⡴⢿⡙⠃⣐⠒⠉⠁⢀⣠⣴⣿⣿⣿⣿⣿⣿⣿⣿⣿⣿⣿⣿⣿⣿⣿⣿⣿⣿⣷⣦⣄⡀⠀⠀⠀⠀⠀⠀⠀⠀⠀⠀⠀⠀⠀⠀⠀⠀⠀⠀⠀⠀⠀⠀⠀⠀⠀⠀⠀⠀⠀⠀⠀⠀
%⣀⣠⣴⣶⣿⣿⣿⣿⣿⣿⣿⣿⣿⣿⣿⣿⣿⣿⣿⣿⣿⣿⣿⣿⣿⣿⣿⣿⣿⣿⣿⣿⣿⣿⣿⣿⣿⣿⣿⣿⣿⣶⣧⣀⡀⠀⠀⠀⠻⠅⡒⠄⢹⡁⠀⠀⠀⢻⠀⠀⠀⢠⣾⣿⡋⢠⠑⢀⣃⣤⣶⣿⣿⣿⣿⣿⣿⣿⣿⣿⣿⣿⣿⣿⣿⣿⣿⣿⣿⣿⣿⣿⣿⣿⣿⣿⣷⣦⣄⣀⠀⠀⠀⠀⠀⠀⠀⠀⠀⠀⠀⠀⠀⠀⠀⠀⠀⠀⠀⠀⠀⠀⠀⠀⠀⠀⠀⠀

\begin{figure}[htbp]
\centering
\includegraphics[width=0.9\linewidth]{overview.png}
\caption{Overview of the three predefined robot manipulation tasks used in this repository.}
\label{fig:ED305-predefined-tasks}
\end{figure}

% ============================================================
\section{Introduction}
% ============================================================

Many robotic systems rely on manually specified behaviors for tasks such as organizing objects or clearing a table. Although this approach can work well in known settings, it may require substantial modification when the objects, environment, or task conditions change. Robot learning instead trains a policy from demonstrations, allowing robots to learn task behaviors without manually programming every possible situation.

The success of a robot learning system depends not only on how it learns, but also on the quality and variety of the data it uses. A robot may perform well in situations it has encountered during training but struggle when faced with different object positions or task settings. Therefore, collecting diverse data and evaluating robot performance under different conditions are important parts of building reliable robot learning systems.

This course project, \textbf{Robot Learning from Evaluation Feedback}, introduces students to an end-to-end process of building and evaluating a robot learning system. Using a simulated NYCU (National Yang Ming Chiao Tung University) AI \& Interaction Specialty Classroom built with NVIDIA Isaac Sim \cite{isaacsim2026} and Isaac Lab \cite{isaaclab} (Figure~\ref{fig:ED305-predefined-tasks}), students will collect data, train and test a robotic arm to perform predefined pick-and-place tasks, and use an ontology-based tool to check whether their data covers the required task conditions. They will also analyze task failures and use the results to improve their data through repeated experiments.

% Simulation provides a safe and repeatable environment for testing different objects, task settings, and robot behaviors. Through this project, students will learn not only how to train a robot to perform a task, but also how to collect good data, evaluate the trained policy, and improve the overall learning process.

After completing the project, students should be able to:
\begin{enumerate}
  \item Understand and carry out an end-to-end robot learning pipeline, from data collection and policy training to evaluation and iterative improvement.
  \item Analyze policy failures, identify possible causes, and iteratively improve the collected dataset through evidence-based adjustments and repeated experiments.
  \item Investigate how demonstration quality, diversity, and task conditions affect policy performance and generalization through systematic experiments and data verification.
\end{enumerate}


% ============================================================
\section{Project Structure}
% ============================================================
This course project consists of \textbf{two levels}: a mandatory \textbf{Entry Level} and an optional \textbf{Advanced Level}. All students across partner institutions are required to complete the Entry Level, which focuses on developing and evaluating an end-to-end robot learning pipeline using predefined manipulation tasks and a shared evaluation framework. The Advanced Level extends this experience by allowing students to define their own manipulation tasks, configure simulation environments, and design appropriate evaluation metrics. Each partner institution determines whether to offer the Advanced Level based on its course arrangements and available computational resources. Students should refer to announcements from their own institution for further details.


\subsection{Entry Level (Mandatory)}

Entry Level standardized setup provides students with hands-on experience in the end-to-end development of Physical AI systems, with an emphasis on the perception-action loop, policy learning, and systematic evaluation. Within the predefined task settings and evaluation framework, each group has the flexibility to explore and modify the following aspects:

\begin{itemize}
  \item \textbf{Training Data:} Data collection methods, environment configs, and dataset size.
  \item \textbf{Policy Design and Deployment:} Policy architectures, training configs, and rollout parameters.
\end{itemize}

Students are encouraged to investigate how these design choices influence policy performance.

\paragraph{Provided Resources.}

The following resources are provided to support students in completing the Entry Level project:
\begin{itemize}
  \item A base simulation environment built on NVIDIA Isaac Sim \cite{isaacsim2026}.
   \item Three predefined pick-and-place tasks (Cup Stacking, Cutlery Arrangement, and Toy Blocks Collection), each provided as a ready-to-use simulation environment with its scene, objects, randomized object start positions, success check, and evaluation configuration (Section~\ref{sec:task-config}).
  \item Two data-collection tools: a state machine controller for each task and keyboard teleoperation (Section~\ref{sec:data-collection}).
  \item Reference policies and learning tools from the LeRobot ecosystem \cite{cadene2024lerobot}.
  \item Standardized evaluation protocols and metrics, such as task success rate.
\end{itemize}

\paragraph{Requirements.}
\label{sec:entry-requirements}

\begin{enumerate}
    \item Select one of the three predefined tasks. Students should refer to announcements from their own institution for task selection.
    \item Configure the provided environment and collect demonstrations for the selected task.
    \item Train and evaluate their policy using the provided tools and evaluation framework.
    \item Iteratively improve their demonstration data quality using the provided ontology-based verification tools, verifying the dataset on at least one of the three generalization axes (Section~\ref{sec:coverage-concept}) and explaining in the report how the requirement was chosen.
\end{enumerate}

\begin{tutornote}
\label{ttn:entry-grading}
    The following table summarizes the grading criteria and their respective weights used by NYCU to assess student performance.
    The Entry Level accounts for 80\% of the final grade and the Advanced Level for the remaining 20\%. The leaderboard
    bonus is added on top.
    \paragraph{Grading Policy.}
    \begin{longtable}{L{0.5\linewidth}L{0.4\linewidth}}
    \toprule
    \textbf{Component} & \textbf{Weight (in final grade)}\\
    \midrule
    Entry Level Technical Report & 40\% \\
    Entry Level Poster Presentation & 40\% \\
    \midrule
    Entry Level total & 80\% \\
    \midrule
    Best private leaderboard success rate in each of the three predefined tasks & +5\% (for the top group of each task) \\
    \bottomrule
    \end{longtable}
\end{tutornote}


\paragraph{Leaderboard.}
\label{sec:leaderboard}

Whether a leaderboard is offered is decided by each institution. When available, groups working on the same task are ranked based on the performance of their policies. The final ranking uses a separate private test configuration to assess how well the policies generalize to unseen scenarios and ensure a fair comparison.

\begin{tutornote}
The link (\url{https://github.com/asddzxcc1856/aicapstone-worker}) describes how to set up leaderboard. Refer to it if needed.
\end{tutornote}


\subsection{Advanced Level (Optional)} 

Compared with the Entry Level, each group has greater flexibility to explore and design the following aspects:

\begin{itemize}
  \item \textbf{Task Definition and Environment Design:} Task objectives, scenes, and success criteria.
  \item \textbf{Training Data:} Data collection methods, dataset size, number of cameras, and camera viewpoints.
  \item \textbf{Policy Design and Deployment:} Policy architectures, training configs, and rollout parameters.
  \item \textbf{Evaluation Design:} Evaluation protocols and performance metrics.
\end{itemize}

Beyond system implementation, students are expected to demonstrate originality in task formulation and system development, justify their design choices, and systematically evaluate their systems. Students are also encouraged to analyze failure cases, identify potential improvements, and consider the challenges of real-world deployment.

\paragraph{Resources.}

Students may draw on publicly available resources, including 3D assets and related implementations from X-Humanoid \cite{artvip_articulated_objects}, Synthesis \cite{synthesis}, and the USD Assets Working Group \cite{usd_assets_working_group} to support their project extensions. \textbf{All Advanced Level work must be conducted in simulation, and students must complete the Entry Level before proceeding to the Advanced Level.}

\paragraph{Requirements.}
\label{sec:advanced-requirements}

\begin{enumerate}
    \item Define an original task with clearly specified objectives, success criteria, and evaluation conditions.
    \item Implement the corresponding simulation environment, including the objects, initial conditions, and task-specific success checks.
    \item Adapt the robot learning pipeline to the newly defined task.
    \item Extend the provided ontology-based data verification tools to support the newly defined task and verify the collected dataset.
    \item Provide complete source code with clear documentation and setup instructions, along with a detailed evaluation protocol specifying random seeds, dataset splits, and other relevant experimental settings.
    \item Provide a reproducibility checklist detailing the procedures required to independently reproduce the experiments and verify the reported results.
\end{enumerate}

\begin{tutornote}
\label{ttn:advanced-grading}
    The following table summarizes the grading criteria and their respective weights used by NYCU to assess student performance.
    \paragraph{Grading Policy.}
    \begin{longtable}{L{0.5\linewidth}L{0.4\linewidth}}
    \toprule
    \textbf{Component} & \textbf{Weight (in final grade)}\\
    \midrule
    Advanced Level Technical Report & 10\% \\
    Advanced Level Poster Presentation & 10\% \\
    \midrule
    Advanced Level total & 20\% \\
    \bottomrule
    \end{longtable}
\end{tutornote}


% ============================================================
\section{Workflow}
\label{sec:structure-workflow}
% ============================================================

\begin{figure}[htbp]
\centering
\includegraphics[width=0.9\linewidth]{flowchart.png}
\caption{Overview of the robot learning workflow, illustrating the iterative feedback loop between data collection, policy training and evaluation, and data verification for continuous dataset improvement.}
\label{fig:project-workflow}
\end{figure}

The workflow of this course project consists of four main stages: task configuration, data collection, policy training and evaluation, and data verification. Figure~\ref{fig:project-workflow} illustrates the overall workflow and the relationships between these stages.

Data verification supports both pre-training dataset assessment and post-evaluation failure analysis. Evaluation and verification results inform subsequent design decisions, enabling an iterative process of system development, evaluation, and refinement.


\subsection{Software Frameworks}
\label{sec:layout-concept}

This course project integrates several open-source software frameworks to support robot simulation, data collection, policy training, and evaluation.

\paragraph{Isaac Sim.}
NVIDIA Isaac Sim \cite{isaacsim2026} is a high-fidelity robotics simulation platform built on NVIDIA Omniverse \cite{nvidia_omniverse_website}. It integrates GPU-accelerated PhysX physics simulation, RTX-based rendering, sensor simulation, and Universal Scene Description (USD), enabling realistic modeling of robot dynamics, environments, and sensors. In this project, Isaac Sim renders the simulated NYCU AI \& Interaction Specialty Classroom (Figure~\ref{fig:ED305-predefined-tasks}) and simulates the Franka robot and the task objects.

\paragraph{Isaac Lab.} Isaac Lab \cite{isaaclab,mittal2025isaaclab} is a modular robot learning framework built on NVIDIA Isaac Sim. It offers standardized interfaces for defining observations, actions, and managing simulation environments. An environment is configured through a collection of \emph{manager} configurations that specify the scene (including the robot, cameras, and objects), observation and action spaces, events (e.g., randomization during reset), termination conditions, and data recorders.

\paragraph{LeIsaac.}
LeIsaac \cite{leisaac} builds on Isaac Lab and integrates with the LeRobot ecosystem \cite{cadene2024lerobot}. It supports teleoperation, automated data collection through state-machine-based scripted policies, and conversion of HDF5 data into the LeRobot dataset format \cite{lerobot_dataset_v3}. This integration enables students to train or fine-tune and deploy various robot learning policies, including Diffusion Policy \cite{chi2023diffusionpolicy}, Action Chunking Transformer (ACT) \cite{zhao2023act}, and NVIDIA Isaac GR00T \cite{gr00tn1_2025}, among others.

\paragraph{Repository Packages.}
Built on top of LeIsaac, the course project repository provides a scripted state-machine controller for each predefined task and keyboard teleoperation for data collection, both recording through the repository's extension of LeIsaac's data recorder. This recorder captures object positions, lighting, object colors, and the task text of each episode for dataset coverage verification. The repository also provides tools for policy rollout, performance evaluation, and data verification.


\subsection{Task Configuration}
\label{sec:task-config}
A task configuration defines the manipulation objective, the simulation environment, the robot and task objects, the initial conditions, and the success criterion. For the Entry Level, three predefined pick-and-place tasks are provided: Cup Stacking, Cutlery Arrangement, and Toy Blocks Collection (Figure~\ref{fig:ED305-predefined-tasks}). Each task randomizes the object positions, the scene lighting, and the wording of the task text at every reset. Section~\ref{sec:task-settings} lists the objects, task texts, and success criteria of each task. Students at the Advanced Level may use the three tasks as examples and starting points for their own tasks.


\subsection{Data Collection}
\label{sec:data-collection-concept}

Data collection is a critical foundation of imitation learning \cite{wiki_imitation_learning}, because the collected 
demonstrations are the only examples from which the policy learns. Two common approaches can be used to collect demonstrations:

\begin{itemize}
\item \textbf{Teleoperation:} A human operator directly controls the robot to demonstrate how a task should be performed.
      This provides flexible and natural demonstrations but requires significant human effort \cite{wikipediaTeleoperation}.
% \item \textbf{Kinesthetic Teaching:} A human physically guides the robot through a task while its movements are recorded.
%       This provides an intuitive way to demonstrate robotic manipulation tasks \cite{botasysKinestheticTeaching}.
\item \textbf{Scripted Methods:} Demonstrations are generated automatically using predefined procedures,
      such as state machines \cite{mdnStateMachine} or motion planners \cite{wikipediaMotionPlanning}. These approaches can
      improve scalability and reduce human effort, but may produce less behavioral diversity.
\end{itemize}

How many episodes to collect, and which conditions they must cover (Section~\ref{sec:coverage-concept}), is a decision that each group makes and justifies. Section~\ref{sec:data-collection} describes the provided collection methods and the recorder.


\subsection{Policy Training and Evaluation}
\label{sec:policies-concept}

% \subsubsection{Imitation Learning}
% Imitation Learning (IL) \cite{wiki_imitation_learning} is a machine learning paradigm in which an agent learns to perform
% tasks by mimicking expert behavior. Unlike Reinforcement Learning (RL) \cite{wiki_reinforcement_learning}, which relies on
% trial and error and a hand-crafted reward function, imitation learning shifts the challenge from designing rewards to
% learning from expert demonstrations.
% 
% The form of imitation learning used in this project is behavior cloning: the policy is trained with supervised
% learning to map observations (the front and wrist camera images and the robot's joint state) to the actions that the
% data-collection controller produced in the same situation. During evaluation, small prediction errors move the robot
% into states that the demonstrations never visited, and errors can compound from there. The performance of an
% imitation-learning policy is therefore strongly influenced by the quality, diversity, and coverage of the
% demonstrations, which is the motivation for data verification
% (Section~\ref{sec:coverage-concept}).

Using the dataset collected in the previous step, students train policies to map observations to robot actions. This repository leverages a suite of robot learning algorithms provided by the LeRobot ecosystem \cite{lerobot_ecosystem}. \textbf{Diffusion Policy} \cite{chi2023diffusionpolicy}, a state-of-the-art imitation learning \cite{wiki_imitation_learning} approach for visuomotor control, is recommended as a starting point. Given visual observations and robot states, the policy learns to predict action sequences (i.e., robot trajectories). Students are encouraged to:
\begin{itemize}
  \item Explore alternative policy architectures within or beyond the provided framework.
  \item Tune training strategies (e.g., data augmentation, normalization, horizon length, etc.) and analyze the relationship between dataset quality and policy performance.
\end{itemize}
This step emphasizes learning-based control, bridging perception and action through a scalable policy learning framework.

After training, the learned policy is executed in the simulator. It receives observations from the simulated environment and predicts actions to control the robot at each simulation step. The resulting execution is recorded and used to evaluate the policy's task performance and generalization ability under different initial conditions. This provides a way to assess whether the learned policy can successfully perform the task beyond the training demonstrations.

Section~\ref{sec:rollout-inputs} lists the observation and action features that rollout expects, and Sections~\ref{sec:training} and~\ref{sec:rollout} describe the training and rollout commands.


\subsection{Data Verification}
\label{sec:coverage-concept}

A failed evaluation does not necessarily mean that the policy architecture or the training procedure is at fault.
Robot learning is highly sensitive to dataset quality, and extensive research has demonstrated its critical role in
building robust systems. For example, \emph{Data Quality in Imitation Learning} demonstrates the importance of data
diversity \cite{belkhale2023data}. RINSE \cite{kulkarni2025learning} and PSD \cite{sojib2026efficient} analyze robot
trajectories in the frequency domain and identify factors that influence policy success rates. More recent works, such
as \emph{Robot Data Curation with Mutual Information Estimators} \cite{hejna2025robot} and DataMiL \cite{dass2026datamil},
quantify data diversity and quality in order to identify the most informative demonstrations.

The course draws on STAR-Gen (A Taxonomy for Evaluating Generalist Robot Manipulation Policies) \cite{gao2026taxonomy},
which organizes the conditions under which a manipulation policy should generalize, such as object position, lighting,
viewpoint, and object appearance (Figure~\ref{fig:star-gen-taxonomy}). The taxonomy is used as a framework for choosing which data factors to verify, not as a definitive standard. Its completeness and long-term relevance remain open questions as policy architectures and evaluation settings evolve, and students are encouraged to examine it critically and to identify both its strengths and its limitations.

\begin{figure}[htbp]
    \centering
    \includegraphics[width=0.9\linewidth]{star-gen.png}
    \caption{STAR-Gen taxonomy for evaluating generalist robot manipulation policies. Reproduced from \cite{gao2026taxonomy}.}
    \label{fig:star-gen-taxonomy}
\end{figure}


The repository checks three axes, each with one factor: Object Poses (VB-POSE; the initial object positions), Image Augmentations (V-AUG; the lighting), and Object Properties (S-PROP; referring to objects by color in the task text). Each group states its requirement for every axis it checks in a \emph{spec} file.

\paragraph{Object Position Coverage.}
\label{sec:object-position-coverage-concept}
As an example, consider the Object Poses axis, which focuses on the initial positions of task-relevant objects. Object positions are randomized on every reset in all three tasks, and the evaluation episodes also start from randomized positions within an area of the table. Object rotation is kept fixed: a random start rotation can leave an object tipped over or make it fall off the table, so the environment has to reset and try again, which wastes simulation time and compute.

If a part of that region has few or no demonstrations, the policy has little to imitate there. The check divides an area of the table into a grid of cells, counts the episodes whose initial object position falls in each
cell, and reports the cells with fewer episodes than required. The area, the cell numbers, and the minimum number of
episodes per cell are chosen by each group because they depend on the task, the evaluation conditions, and the
collection budget.

The check is built from five knowledge-representation tools:

\begin{itemize}
  \item \textbf{Ontology.} An ontology \cite{guarino2009ontology} provides a formal representation of knowledge in a specific domain by defining its shared concepts and the relationships among them. In this repository, the ontology describes the concepts relevant to dataset coverage verification, including generalization axes, datasets, episodes, task objects, factor values such as initial object positions, and coverage requirements.
  \item \textbf{Web Ontology Language (OWL).} OWL \cite{w3c2012owl} is a standard language for representing ontologies. It provides formal semantics based on Description Logic (DL), allowing domain knowledge and constraints to be expressed in a machine-interpretable form.
  \item \textbf{RDF knowledge graph in Turtle.} The recorded dataset, the group's coverage requirement, each rollout,
        and every verification result are expressed as Resource Description Framework (RDF) facts over this vocabulary
        and kept in one \emph{verification store} per dataset, a directory of Turtle files \cite{w3cTurtle}. The
        requirement itself is written as a small Turtle file and attached to the store.
  \item \textbf{SHACL validation.} Shapes Constraint Language (SHACL) shapes \cite{w3cShacl} turn the requirement into executable constraints: they
        compare factor values (for example positions) against the required bounds, count episodes per cell, and report violations. This step is needed
        because Description Logic alone does not perform numeric interval comparison or counting.
  \item \textbf{SPARQL queries.} The store records what was recorded, what was required, and what verification found,
        so students can ask their own questions of it with SPARQL instead of relying only on the
        printed report.
\end{itemize}

The store separates raw capture from interpretation. The LeRobot dataset and the rollout log stay the sources: the
dataset is still uploaded and trained on as before, and the rollout log is written line by line during the rollout,
so an interrupted rollout still leaves a readable log. Both enter the store afterwards, on the host, with their path and a
content hash, so the tools can tell when a source has changed since it entered the store.

Coverage is a diagnostic, not a proof of the cause of a failure. An uncovered evaluation region is evidence of a
possible gap in the data distribution, while complete coverage does not guarantee that a policy will succeed. Failures
should be treated as feedback: inspected in detail, tested with experiments, and used to refine the collection. Section~\ref{sec:data-verification} describes how to run the check.



% ============================================================
\section{Setup}
\label{sec:setup}
% ============================================================

\subsection{Required Resources}

\begin{itemize}
  \item a Linux x86\_64 machine with an NVIDIA RTX GPU with at least 16~GB of memory and driver 570 or newer. Recording and rollout render camera images, which requires a GPU with RT cores (GeForce RTX, RTX A/Ada/PRO, L40S). Data-center GPUs without RT cores (A100, H100, V100) can be used for training only;
  \item Ubuntu 22.04 or newer (the Isaac Sim packages need glibc 2.35 or newer), about 40~GB of free disk space for the simulator environment, and a desktop session for the Isaac Sim window;
  \item Python 3.11 and \code{uv} (the Python version is pinned in \code{.python-version});
  \item Git with submodule support;
  \item a Hugging Face account and access token, for transferring datasets and checkpoints;
  \item a Weights \& Biases account, for monitoring training;
  \item optionally, Docker and the NVIDIA Container Toolkit, only for the container alternative to the host simulator installation (Section~\ref{sec:env-setup}).
\end{itemize}

Check the tools:
\begin{lstlisting}[language=bash]
uv --version
make --version
nvidia-smi
\end{lstlisting}


\subsection{GlowsAI}

Each group receives 30 GlowsAI credits for the course project. Follow the GlowsAI setup guide
(\url{https://docs.google.com/presentation/d/1xPHm3gNNvLbX7aOWnY_xW5dfVLEa9XiBuskMbRsgxI8/edit?usp=sharing})
to create a GPU instance. When you use the terminal of the Virtual Network Computing (VNC) desktop on GlowsAI, switch to the \code{glows} user first:

\begin{lstlisting}[language=bash]
su - glows
\end{lstlisting}


\subsection{Fork and Clone the Repository}
\label{sec:fork}

You do not have write access to the course repository (\url{https://github.com/HCIS-lab/RobotLearningFromEvaluationFeedback}). Each group therefore works in its own \emph{fork}, a copy of the repository under a group member's GitHub account. Your changes (state-machine and task configuration changes, filled-in specs) are committed and pushed to the fork, so they are kept and shared within the group. A plain clone of the course repository would leave all changes on one machine.

\begin{tcolorbox}[colback=yellow!8,colframe=orange!70!black,boxrule=0.4pt,arc=1mm]
\textbf{Never push to the course repository (\code{HCIS-lab/RobotLearningFromEvaluationFeedback}).} Push only to your fork (\code{origin}). The remote named \code{upstream} is read-only for you: use it only to \emph{fetch} course updates.
\end{tcolorbox}

\begin{enumerate}
  \item One group member opens the course repository on GitHub, clicks \textbf{Fork}, and creates the fork under their own account. The other members are added as collaborators (fork \textbf{Settings} $\rightarrow$ \textbf{Collaborators}).
  \item Every member clones the \emph{fork}, not the course repository, and adds the course repository as \code{upstream}:
\end{enumerate}

\begin{lstlisting}[language=bash]
git clone https://github.com/<github-user>/RobotLearningFromEvaluationFeedback   # your fork
cd RobotLearningFromEvaluationFeedback
git remote add upstream https://github.com/HCIS-lab/RobotLearningFromEvaluationFeedback
git remote -v    # origin = your fork (push here), upstream = course repository (fetch only)
\end{lstlisting}

Commit and push your work to the fork:

\begin{lstlisting}[language=bash]
git add <files>
git commit -m "<message>"
git push origin main     # always origin, never upstream
\end{lstlisting}

When the course repository is updated (fixes, new evaluation files), merge the updates into your fork. Conflicts appear only where your group changed the same lines as the update.

\begin{lstlisting}[language=bash]
git fetch upstream
git merge upstream/main
git push origin main
\end{lstlisting}

Datasets, checkpoints, and training outputs (\path{datasets/}, \path{checkpoints/}, \path{outputs/}) are ignored by Git. They are shared through the Hugging Face Hub (Section~\ref{sec:hf}), not through the fork.


\subsection{Environment Setup}
\label{sec:env-setup}

On the host, in the cloned fork:

\begin{lstlisting}[language=bash]
make submodules                    # Isaac Lab submodule; the launch targets also run it
uv sync                            # host environment for training and the coverage check
source .venv/bin/activate
hf auth login                      # paste your Hugging Face access token
wandb login                        # paste your Weights & Biases API key
export HF_USER=<your-huggingface-username>
\end{lstlisting}

The simulator (Isaac Sim 5.1.0, Isaac Lab, LeIsaac, and LeRobot) is installed into its own environment, \code{.venv-sim}, separate from the \code{uv} environment \code{.venv}, because Isaac Sim needs older package versions (NumPy 1.26, PyTorch 2.7) than training and verification. Install it once on the host (this downloads about 20~GB) and open a simulator shell:

\begin{lstlisting}[language=bash]
make install-sim    # once; downloads about 20 GB, takes a while
make launch-sim     # opens a shell with .venv-sim active, in the repository root
\end{lstlisting}

Type \code{exit} to leave the shell. In an existing terminal, \code{source scripts/sim_env.sh} does the same. The shell accepts the Isaac Sim license and points LeRobot's dataset folder to \path{datasets/}, so a recorded dataset appears in \path{datasets/<HF_USER>/<name>}. Recording and rollout run in this shell. Run \code{hf auth login} and set \code{<HF_USER>} in it as well if you transfer data from it. The first start of Isaac Sim compiles shaders and takes several minutes. The Isaac Sim window opens on the display in \code{DISPLAY}, so start the shell from a terminal on the machine's desktop; commands with \code{--headless} run without a window. Training and verification commands (\code{uv run ...}, \code{lerobot-train}) run outside the simulator shell.

\paragraph{Alternative: Isaac Lab container.} If the host installation is not possible, build and enter the Isaac Lab container instead. Choose the target that matches your machine; the first build takes a while:

\begin{lstlisting}[language=bash]
make launch-isaaclab               # local machine
make launch-isaaclab-glowsai-4090  # RTX 4090
make launch-isaaclab-glowsai-l40s  # L40S
\end{lstlisting}

The command drops you into a shell inside the container, with the repository mounted at
\path{/workspace/physicalai}. All commands in this document work the same in this shell. Run \code{make check-isaaclab-gpu} if Isaac Sim cannot find the GPU, and \code{make fix-perms} to give files that the container wrote under \path{datasets/} and \path{outputs/} back to your host user. Run \code{make test} on the host to confirm that the Python environment works.


\subsection{Hugging Face Upload and Download}
\label{sec:hf}

Before uploading or downloading, make sure that you are logged in:
\begin{lstlisting}[language=bash]
hf auth login    # log in if needed
hf auth whoami   # shows the logged-in user
\end{lstlisting}

Upload a dataset (on the host, recorded datasets are under \path{datasets/}) or a trained checkpoint:
\begin{lstlisting}[language=bash]
hf upload <HF_USER>/<repo_id> datasets/<HF_USER>/<repo_id> --repo-type dataset
hf upload <HF_USER>/<policy_repo> <output_dir>/checkpoints/last/pretrained_model --repo-type model
\end{lstlisting}

Download a dataset or a checkpoint:
\begin{lstlisting}[language=bash]
hf download <HF_USER>/<repo_id> --repo-type dataset --local-dir datasets/<HF_USER>/<repo_id>
hf download <HF_USER>/<repo_id> --repo-type model --local-dir checkpoints/<repo_id>
\end{lstlisting}

Add \code{--revision <branch, tag, or commit>} to download a specific version. Record the commit hash of every dataset
and checkpoint that you report, so that your results can be reproduced.


% ============================================================
\section{Implementation Details}
\label{sec:implementation}
% ============================================================

This section guides students through using the repository after covering the concepts in Section~\ref{sec:structure-workflow} and the setup in Section~\ref{sec:setup}. Guides with more detail are in \path{docs/}.

\subsection{Execution Environments}
\label{sec:exec-env}

The repository uses two independent Python environments on the same machine (Section~\ref{sec:env-setup}).

\begin{longtable}{L{0.25\linewidth}L{0.30\linewidth}L{0.37\linewidth}}
\toprule
\textbf{Stage} & \textbf{Environment} & \textbf{Reason} \\
\midrule
Data collection, policy rollout & Simulator shell (\path{.venv-sim}) or the \code{isaaclab} container & Both stages need Isaac Sim. \\
Policy training & \code{uv} environment (\path{.venv}) & The simulator is not needed. \\
Data verification & \code{uv} environment (\path{.venv}) & It only reads files and does not require the simulator. \\
\bottomrule
\end{longtable}


\subsection{Repository Overview}
\label{sec:repo-overview}

\begin{longtable}{L{0.36\linewidth}L{0.58\linewidth}}
\toprule
\textbf{Path} & \textbf{Content} \\
\midrule
\endhead
\path{packages/simulator/src/simulator/tasks/template/single_arm_franka_cfg.py}
  & Shared task template: Franka robot, front and wrist cameras, observations, action space, and the LeRobot frame builder. \\
\path{packages/simulator/src/simulator/tasks/<task>/}
  & One directory per task. \code{<task>_env_cfg.py} defines the scene, objects, randomization events, and \code{success} termination term; \code{__init__.py} registers the Gym ID. \\
\path{packages/simulator/src/simulator/datagen/state_machine/}
  & One scripted state-machine controller per task. \\
\path{packages/simulator/src/simulator/devices/}
  & Keyboard device for teleoperation. \\
\path{packages/simulator/src/simulator/utils/}
  & \code{lerobot_recorder.py}: the repository's \code{LeRobotRecorderManager}. \code{domain_randomization.py}: reset events (object positions, lighting). \code{rollout_log.py}: writer of the rollout log. \\
\path{packages/simulator/assets/}
  & USD scenes and object assets. \\
\path{packages/data_verification/}
  & Verification store: dataset and rollout-log readers, OWL ontology, SHACL rules, and one verifier per generalization axis under \path{axes/}. \\
\path{eval/}
  & Evaluation file of each task, used by \code{rollout.py}. \\
\path{specs/<task>/}
  & Shipped spec of each task, one per axis, with every value set to \code{TODO}. \\
\path{queries/}
  & Example SPARQL queries. \\
\path{scripts/datagen/generate.py}, \path{scripts/teleop.py}
  & Data collection with the state machine and with keyboard teleoperation. \\
\path{scripts/store.py}
  & Command-line interface of the verification store (\code{ingest}, \code{verify}, \code{explore}, \code{query}). \\
\path{scripts/rollout.py}
  & Runs a trained LeRobot checkpoint in the simulator and reports the success rate. \\
\path{docs/}, \path{tests/}
  & Markdown guides, and unit tests (\code{make test}). \\
\path{Dockerfile}, \path{Makefile}, \path{requirements-sim.txt}
  & Container image, the \code{make} targets, and the pinned versions of the simulator environment. \\
\path{datasets/}, \path{outputs/}, \path{checkpoints/}, \path{stores/}
  & Recorded datasets, training and rollout outputs, checkpoints, and verification stores. Not tracked by Git. \\
\bottomrule
\end{longtable}


\subsection{Tasks}
\label{sec:task-settings}

The three predefined tasks are summarized below. Figure~\ref{fig:task-overview} shows them.

\begin{longtable}{L{0.17\linewidth}L{0.24\linewidth}L{0.24\linewidth}L{0.24\linewidth}}
\toprule
\textbf{Task} & \textbf{Cup Stacking} & \textbf{Cutlery Arrangement} & \textbf{Toy Blocks Collection} \\
\midrule
\endhead
Task ID & \code{HCIS-CupStacking-SingleArm-v0} & \code{HCIS-CutleryArrangement-SingleArm-v0} & \code{HCIS-ToyBlocksCollection-SingleArm-v0} \\
Objects in \code{env.scene} & \code{blue_cup}, \code{pink_cup} & \code{plate}, \code{fork}, \code{knife} & \code{green_block}, \code{blue_block}, \code{red_block}, \code{storage_box} \\
Task text, by name & ``pick up the blue cup and place it on the pink cup.'' & ``place the fork on the left and the knife on the right of the plate.'' & ``pick up the bridge, the cylinder and the triangle and place them into the storage box.'' \\
Task text, by color & ``pick up the blue object and place it on the pink object.'' & ``place the pink object on the left and the orange object on the right of the plate.'' & ``pick up the green object, the blue object and the red object and place them into the storage box.'' \\
Success term & Blue cup within 2~cm of the pink cup in x and in y, and 8--10~cm higher. & Fork and knife each inside a $10\times10$~cm square centred 12~cm from the plate centre; fork at smaller x (the robot's left), knife at larger x. & Every block inside the storage-box footprint ($0.296\times0.209$~m) shrunk by 1~cm per side (box frame), and 0--11~cm above the box bottom. \\
Episode length in rollout & 20~s & 40~s & 60~s \\
\bottomrule
\end{longtable}

\begin{figure}[htbp]
    \centering
    \begin{subfigure}[b]{0.32\textwidth}
        \centering
        \includegraphics[width=\textwidth]{cs_overview.png}
        \caption{Cup Stacking}
        \label{fig:cs_overview}
    \end{subfigure}
    \hfill
    \begin{subfigure}[b]{0.32\textwidth}
        \centering
        \includegraphics[width=\textwidth]{ca_overview.png}
        \caption{Cutlery Arrangement}
        \label{fig:ca_overview}
    \end{subfigure}
    \hfill
    \begin{subfigure}[b]{0.32\textwidth}
        \centering
        \includegraphics[width=\textwidth]{tbc_overview.png}
        \caption{Toy Blocks Collection}
        \label{fig:tbc_overview}
    \end{subfigure}
    \caption{Overview of the three predefined tasks.}
    \label{fig:task-overview}
\end{figure}

The randomization at every reset (Section~\ref{sec:task-config}) shifts the object positions by a few centimetres, changes the intensity and color temperature of the scene light \code{rect_light}, and picks the name or color form of the task text. Object rotation and color never change.

The task configurations are in \path{packages/simulator/src/simulator/tasks/<task>/<task>_env_cfg.py}. The shared robot, camera, action, and timing settings are in \path{template/single_arm_franka_cfg.py}.


\subsection{Data Collection}
\label{sec:data-collection}

Record demonstrations in the simulator shell with one of two provided methods. Every recording command needs \code{--enable_cameras} because the dataset contains camera images.

\begin{longtable}{L{0.18\linewidth}L{0.38\linewidth}L{0.38\linewidth}}
\toprule
& \textbf{State machine} & \textbf{Keyboard teleoperation} \\
\midrule
Script & \path{scripts/datagen/generate.py} & \path{scripts/teleop.py} \\
Controller & A scripted state machine per task & A person, through the keyboard \\
Saved episodes & Those that pass the state machine's own success check & Those that the operator marks with success \\
Display & GUI or \code{--headless} & GUI only \\
\bottomrule
\end{longtable}

\subsubsection{State Machine}
\label{sec:state-machine}

\begin{lstlisting}[language=bash]
python scripts/datagen/generate.py \
    --task <task_id> \
    --num_demos 50 \
    --device cuda --enable_cameras --headless \
    --record --use_lerobot_recorder \
    --lerobot_dataset_repo_id <HF_USER>/<repo_id>
\end{lstlisting}

\begin{itemize}
  \item \code{--task} specifies one of the following task IDs: \code{HCIS-CupStacking-SingleArm-v0}, \code{HCIS-CutleryArrangement-SingleArm-v0}, or \code{HCIS-ToyBlocksCollection-SingleArm-v0}.
  \item \code{--num_demos} is the number of attempts, successful or not. Only successful attempts are saved, so the dataset usually has fewer episodes.
  \item \code{--record --use_lerobot_recorder --lerobot_dataset_repo_id <HF_USER>/<repo_id>} records into \path{datasets/<HF_USER>/<repo_id>}. Without \code{--record}, nothing is saved.
  \item \code{--resume} appends to an existing dataset with the same id. \code{--seed} sets the random seed (default: current time).
\end{itemize}

The state machine reads the object positions after each reset, moves through its phases (for example hover, approach, grasp, lift, move, lower, release), and aborts the attempt if an object falls below the table.

\subsubsection{Keyboard Teleoperation}
\label{sec:teleop}

\begin{lstlisting}[language=bash]
python scripts/teleop.py \
    --task <task_id> \
    --teleop_device keyboard \
    --device cuda --enable_cameras \
    --record --use_lerobot_recorder \
    --lerobot_dataset_repo_id <HF_USER>/<repo_id> \
    --num_demos 20
\end{lstlisting}

\code{--num_demos} stops the session after that many saved episodes (0, the default, runs until Ctrl+C). \code{--sensitivity} scales the motion per key press. The other options are as for the state machine. The Isaac Sim window needs the keyboard focus. The keys move a target pose of the hand, and an inverse kinematics solver turns it into joint targets:

\begin{longtable}{L{0.30\linewidth}L{0.62\linewidth}}
\toprule
\textbf{Key} & \textbf{Action} \\
\midrule
B & Start control; recording of the episode starts here \\
W / S, A / D, K / J & Move the hand forward (+y) / backward ($-$y), left ($-$x) / right (+x), and up (+z) / down ($-$z), in the Isaac Sim world frame. \\
H / L & Pitch ($-$ / $+$) \\
U / I & Roll ($-$ / $+$) \\
E / Q & Yaw ($-$ / $+$) \\
C / M & Open / close the gripper \\
N & Save the episode and reset \\
R & Discard the episode and reset \\
\bottomrule
\end{longtable}

The script does not apply the task's success check. Press \textbf{N} only when the task is solved as defined in Section~\ref{sec:task-settings}, and \textbf{R} otherwise. After every reset, press \textbf{B} to start control again.

\subsubsection{Adjusting the Collection}
\label{sec:adjusting}

Groups may adjust the collection to target the data they need, and report the changes as part of their data-collection design. For example:

\begin{itemize}
  \item the state machine constants at the top of \path{datagen/state_machine/<task>.py}.
  \item the object randomization ranges, start positions, or success check in the task configuration.
  \item the light ranges \code{intensity_range} and \code{color_temperature_range} of the \code{randomize_light_conditions} event.
  \item how often the task text refers to the objects by color, with \code{reference_weights}.
\end{itemize}

Evaluation always uses the files in \path{eval/}, so these changes affect only the recorded data.


\subsubsection{Recorder and Dataset}
\label{sec:recorder}

Students are required to record with \code{--use_lerobot_recorder}, which attaches the repository's \code{LeRobotRecorderManager} and writes a LeRobot v3 dataset. Without it, \code{--record} writes an HDF5 file that can be neither trained on nor verified. It extends LeIsaac's recorder with the columns that data verification reads, in addition to the rollout features of Section~\ref{sec:rollout-inputs}:

\begin{longtable}{L{0.30\linewidth}L{0.62\linewidth}}
\toprule
\textbf{Column} & \textbf{Content} \\
\midrule
\endhead
\code{object_position.<name>} & \code{(3,)} float32. Position $[x,y,z]$ in metres, relative to the environment origin, of each tracked object. Fixed objects such as \code{plate} and \code{storage_box} have no column. \\
\code{lighting.rect_light} & \code{(2,)} float32. Light intensity and color temperature in kelvin. \\
\code{object_color.<name>} & \code{(3,)} float32. Diffuse color $[r,g,b]$ of each object that the task text refers to. It shows that the object itself never changes. \\
\code{task_index} & int64. Index of the episode's task text in the dataset's \path{meta/tasks.parquet} file. \\
\bottomrule
\end{longtable}

Language-conditioned policies (for example GR00T, $\pi_0$, or SmolVLA) use the task text as instruction, whereas Diffusion Policy and ACT ignore it. 

Further behavior of the recorder:

\begin{itemize}
  \item One frame is written per \code{env.step()} at 60~fps. The first 5 steps after each reset are skipped.
  \item An episode is saved at the next reset if it succeeded, and dropped otherwise.
  \item The dataset goes to \path{datasets/<HF_USER>/<repo_id>}. If the folder exists, recording stops with \code{FileExistsError}; use \code{--resume} to append.
  \item Stop with Ctrl+C. A process killed with \code{kill -9} leaves a dataset that cannot be loaded.
\end{itemize}

\subsubsection{Check Your Dataset}
\label{sec:check-dataset}

\begin{enumerate}
  \item \textbf{Inspect the collected dataset} using the LeRobot dataset visualizer \cite{lerobotVisualizer}, which loads datasets by their Hugging Face repository ID. Verify that the first frame shows the reset scene, the last frame shows task completion, both cameras display the correct views, and the state curves reflect the robot's movements.
  \item \textbf{Load it and look at the numbers:}
    \begin{lstlisting}[language=Python]
    from lerobot.datasets.lerobot_dataset import LeRobotDataset

    ds = LeRobotDataset("<HF_USER>/<repo_id>", root="datasets/<HF_USER>/<repo_id>")
    print(ds.num_episodes, ds.num_frames, ds.fps)
    print(list(ds.features))
    \end{lstlisting}
  \item \textbf{Verify the dataset against your specs} (Section~\ref{sec:data-verification}). It is a coverage diagnostic and does not replace visual inspection.
\end{enumerate}


\subsection{Data Verification}
\label{sec:data-verification}

All verification goes through one \emph{verification store} per dataset (\path{stores/<repo_id>/}), which holds the dataset, your specs, the rollouts, and every result as Turtle files. The commands run in the \code{uv} environment from the repository root. Each axis has a guide in \path{docs/}: \path{object_pose.md}, \path{image_augmentation.md}, and \path{object_property.md}.

\subsubsection{Ingest the Dataset}

\begin{lstlisting}[language=bash]
uv run python scripts/store.py ingest --dataset datasets/<HF_USER>/<repo_id>
\end{lstlisting}

This creates \path{stores/<repo_id>/dataset.ttl}: each episode with its frame count, task text, and the value of every factor at its first recorded frame. Only the metadata and Parquet files are read, so it takes seconds. Run it again after recording more episodes; it removes the previous results, which described the old episodes.

\subsubsection{Write the Specs}
\label{sec:specs}

A spec is a small Turtle file that states your requirement for one axis. The repository ships one per axis and task in \path{specs/<task>/}, with every value set to \code{"TODO"}:

\begin{longtable}{L{0.32\linewidth}L{0.25\linewidth}L{0.30\linewidth}}
\toprule
\textbf{File} & \textbf{Axis (STAR-Gen)} & \textbf{Factor checked} \\
\midrule
\code{object_pose.ttl} & Object Poses (VB-POSE) & object position \\
\code{image_augmentation.ttl} & Image Augmentations (V-AUG) & lighting \\
\code{object_property.ttl} & Object Properties (S-PROP) & referencing objects by color \\
\bottomrule
\end{longtable}

Copy the files of the axes you are going to check, for example into \path{my_specs/}, and replace every \code{"TODO"}. Numbers are written without quotes; text keeps its quotes. Each file is commented, and the axis guide explains how to choose each value. An axis without a spec file is not checked. The file name is the spec's name in reports.

\subsubsection{Verify and Read the Report}
\label{sec:verify}

\begin{lstlisting}[language=bash]
uv run python scripts/store.py verify --store stores/<repo_id> \
    --spec my_specs/object_pose.ttl my_specs/object_property.ttl
\end{lstlisting}

\code{verify} checks each spec file, checks the dataset against it, prints one report per spec, and writes the results to \path{stores/<repo_id>/findings/}. Each report starts with the store, the spec, and the number of episodes, followed by one block per requirement (see the axis guides). The lines are:

\begin{itemize}
  \item \textbf{ERROR}: the check could not run, for example a \code{TODO} is left or a column is missing.
  \item \textbf{VIOLATION}: a requirement is not met, for example a cell has too few episodes.
  \item \textbf{WARNING}: allowed but worth knowing, for example an episode outside the required area. It does not count toward coverage.
  \item \textbf{RESULT}: \code{coverage spec met} when the spec has no error and no violation.
\end{itemize}

The exit code is 0 when every spec is met and 1 otherwise. Useful options: \code{--strict} also fails on warnings, \code{--no-rollouts} skips rollout comparison, and \code{--json <file>} writes the full result. The store keeps exactly the specs of the last \code{verify}. \code{explore} summarizes the store:

\begin{lstlisting}[language=bash]
uv run python scripts/store.py explore --store stores/<repo_id>
\end{lstlisting}

\subsubsection{Query the Store}

Every fact and result in the store can be queried with SPARQL, locally and without a server. Examples are in \path{queries/}, for instance \code{uncovered_cells.rq}, which lists the cells that need more episodes:

\begin{lstlisting}[language=bash]
uv run python scripts/store.py query --store stores/<repo_id> \
    --query-file queries/uncovered_cells.rq
uv run python scripts/store.py query --store stores/<repo_id> \
    --query-text 'SELECT (COUNT(?e) AS ?n) WHERE { ?e a cp:Episode }'
\end{lstlisting}

\code{SELECT} prints a table, and the prefix \code{cp:} is predefined. Compare numbers with \code{FILTER}, because indices are stored as \code{xsd:nonNegativeInteger}.


\subsection{Training}
\label{sec:training}

Training runs in the \code{uv} environment (Section~\ref{sec:exec-env}).
\begin{lstlisting}[language=bash]
lerobot-train \
  --dataset.repo_id=<HF_USER>/<repo_id> \
  --policy.type=diffusion \
  --policy.device=cuda \
  --policy.push_to_hub=false \
  --policy.n_obs_steps=2 --policy.horizon=16 --policy.n_action_steps=8 \
  --policy.optimizer_lr=1e-4 \
  --output_dir=outputs/train/<job_name> \
  --job_name=<job_name> \
  --steps=100000 --save_freq=20000 --log_freq=1000 \
  --wandb.enable=true
\end{lstlisting}

\begin{itemize}
  \item \code{--dataset.repo_id} downloads the dataset from the Hugging Face Hub. To train on a local copy, also pass
        \code{--dataset.root=datasets/<HF_USER>/<repo_id>}.
  \item \code{--policy.type} selects the policy \cite{cadene2024lerobot}, e.g.\ \code{groot}, \code{diffusion}, or \code{act}. The \code{--policy.*} options after it are policy-specific; the line above uses Diffusion Policy options. For ACT, use options such as \code{--policy.chunk_size} and \code{--policy.n_action_steps}.
  \item \code{--policy.n_action_steps} is the number of actions that rollout executes open loop after each inference,
        each lasting one 1/60~s step. \code{rollout.py} uses the value stored in the checkpoint unless you pass
        \code{--policy_action_horizon}.
  \item Reduce \code{--batch_size} if you run out of GPU memory.
\end{itemize}

Checkpoints are written to \path{<output_dir>/checkpoints/<step>/pretrained_model}, and \path{checkpoints/last} points to
the latest one. Each checkpoint contains \code{config.json}, \code{model.safetensors}, and the pre- and post-processor files
with the normalization statistics. See the LeRobot documentation \cite{cadene2024lerobot} for all options.


\subsection{Rollout}
\label{sec:rollout}

Run the evaluation in the simulator shell:

\begin{lstlisting}[language=bash]
python scripts/rollout.py \
  --task=eval/cup_stacking_eval.py \
  --policy_type=lerobot-diffusion \
  --policy_checkpoint_path=outputs/train/<job_name>/checkpoints/last/pretrained_model \
  --device=cuda --enable_cameras \
  --eval_rounds=10
\end{lstlisting}

\begin{longtable}{L{0.32\linewidth}L{0.62\linewidth}}
\toprule
\textbf{Flag} & \textbf{Meaning} \\
\midrule
\endhead
\code{--task} & The path to an evaluation \code{.py} file: \code{eval/cup_stacking_eval.py}, \code{eval/cutlery_arrangement_eval.py}, or
               \code{eval/toy_blocks_collection_eval.py}. \\
\code{--policy_type} & \code{lerobot-<type>}, matching \code{--policy.type} in training, e.g.\ \code{lerobot-diffusion} or
               \code{lerobot-act}. Required; one of \code{lerobot-diffusion}, \code{-act}, \code{-smolvla}, \code{-pi0},
               \code{-pi05}, \code{-vqbet}, \code{-tdmpc}, \code{-sac}, \code{-groot}, \code{-xvla}. \\
\code{--policy_checkpoint_path} & Local \code{pretrained_model} directory of the checkpoint. \\
\code{--eval_rounds} & Number of evaluation episodes. \\
\code{--policy_action_horizon} & Optional. Actions executed per inference; overrides \code{n_action_steps} of the
               checkpoint. Must be positive and at most the actions one inference predicts. \\
\code{--episode_length_s} & Timeout of each episode in seconds. Default: the task's graded length (20, 40, or 60~s). \\
\code{--seed} & Seed of the object position and lighting randomization (default 42). Keep it fixed to compare checkpoints. \\
\code{--headless} & Run without the GUI, which is faster. \\
\bottomrule
\end{longtable}

An episode ends as soon as the \code{success} term of the evaluation file becomes true, or at the timeout. The script
prints the episode length and action horizon in use at start, the success rate after every round, and the final
success rate at the end.

Every rollout also writes a log, \path{outputs/rollout/<timestamp>.jsonl}, without any flag. It has a header line (evaluation file, seed, checkpoint) and one line per episode with its status, the factor values it started with, and the task text given to the policy.


\paragraph{Rollout Input Features.}
\label{sec:rollout-inputs}

You may customize the input and output features in your experiments. Just remember to keep your policy consistent with the rollout input feature specification. The default input features are:

\begin{longtable}{L{0.31\linewidth}L{0.12\linewidth}L{0.15\linewidth}L{0.37\linewidth}}
\toprule
\textbf{Key} & \textbf{dtype} & \textbf{shape} & \textbf{meaning} \\
\midrule
\code{observation.state} & float32 & \code{(9,)} & Measured joint positions, in the order \code{panda_joint1..7, panda_finger_joint1, panda_finger_joint2}. Arm in rad, fingers in m. \\
\code{action}  & float32 & \code{(8,)} & \code{[0:7]} absolute joint position targets in rad; \code{[7]} gripper: $< 0$ closes, $\geq 0$ opens. Use $\pm 1$. \\
\code{observation.images.front}  & image/video & \code{(480, 640, 3)} & RGB uint8, height $\times$ width $\times$ channels (HWC), from the fixed front camera. \\
\code{observation.images.wrist}  & image/video & \code{(480, 640, 3)} & RGB uint8, HWC, from the wrist camera. \\
\bottomrule
\end{longtable}


\subsection{Compare Your Dataset with a Rollout}
\label{sec:rollout-compare}

The dataset check compares the dataset with your own requirement. After the rollout has ended, add the log to the dataset store used to train the policy, then verify it against the same specifications:

\begin{lstlisting}[language=bash]
uv run python scripts/store.py ingest --store stores/<repo_id> \
    --rollout-log outputs/rollout/<timestamp>.jsonl
uv run python scripts/store.py verify --store stores/<repo_id> \
    --spec my_specs/object_pose.ttl my_specs/object_property.ttl
\end{lstlisting}

For every ingested rollout and every spec, \code{verify} prints a report. For each completed run, each axis independently determines whether the dataset covers it. A run is considered in-distribution only if it is covered across all axes specified in the spec. The report compares the success rate inside and outside the distribution, and suggests what to record next. A much lower success rate outside points to missing data rather than the model. With two or more spec files, a final \code{COMBINED} block counts a run as in distribution only when every spec file does. Repeating record, train, rollout, and verify closes the gaps round by round.


% ============================================================
\section{Deliverables}
% ============================================================

Submission deadlines and specific penalty policies vary by institution. Please follow the announcements and designated submission platform provided by your institution. The required report contents, the files to submit, and their names and directory structure will be given in a separate document, announced to students.

\begin{itemize}
    \item \textbf{Deadline, Late Submission, and Submission Format: }
    Follow the deadline, the late-submission penalty, and the file naming, directory structure, and report format announced by your institution. Incorrect submission formats, including a report not written in English, may result in a penalty according to your institution's policy.

    \item \textbf{Plagiarism: }
    All submitted work must be original. Plagiarism, including unauthorized copying of code or report content from other students or public sources, is prohibited and will be handled according to your institution's policy.

    \item \textbf{Repository Content: }
    Please do not commit sensitive information, such as access tokens, or large runtime artifacts, such as datasets and checkpoints, to the repository. Instead, share datasets and checkpoints through the Hugging Face Hub.
\end{itemize}


\begin{tutornote}
The following table provides NYCU's format requirement as a reference. Grading details please refer to Tutornote~\ref{ttn:entry-grading} and Tutornote~\ref{ttn:advanced-grading}.
\begin{longtable}{p{0.2\linewidth}p{0.7\linewidth}}
\toprule
\textbf{Criterion} & \textbf{Indicators} \\
\midrule
Report & Must include the task description, data collection pipeline, policy training, experimental results, evaluations conducted and resulting improvements, and key learning, discussion, and insights. \\
Poster Presentation & In English, and must include the task description, pipeline and method modifications, evaluations conducted and resulting improvements, and key learning and insights. \\
\bottomrule
\end{longtable}
\end{tutornote}


\printbibliography[heading=bibintoc]
\end{document}